
\documentclass{article}
\usepackage{colortbl}
\usepackage{makecell}
\usepackage{multirow}
\usepackage{supertabular}

\begin{document}

\newcounter{utterance}

\centering \large Interaction Transcript for game `cladder', experiment `full\_v1.5\_default', episode 3752 with llama{-}3.1{-}8b{-}instruct.
\vspace{24pt}

{ \footnotesize  \setcounter{utterance}{1}
\setlength{\tabcolsep}{0pt}
\begin{supertabular}{c@{$\;$}|p{.15\linewidth}@{}p{.15\linewidth}p{.15\linewidth}p{.15\linewidth}p{.15\linewidth}p{.15\linewidth}}
    \# & \multicolumn{2}{c}{Player} && \multicolumn{2}{c}{Game Master} \\
    \hline

    \theutterance \stepcounter{utterance}  
    & & & \multicolumn{4}{p{0.6\linewidth}}{
        \cellcolor[rgb]{0.9,0.9,0.9}{
            \makecell[{{p{\linewidth}}}]{
                \texttt{\tiny{[P1$\langle$GM]}}
                \texttt{You are an expert in causal inference. The following question is not a typical commonsense query, but rather a meticulously designed question created by a professor specializing in causal inference, intended to assess the students' mastery of the course content.} \\
\texttt{Imagine a self{-}contained, hypothetical world with only the following conditions, and without any unmentioned factors or causal relationships: Listening to jazz has a direct effect on elite institution admission status. Effort has a direct effect on elite institution admission status. For students accepted to elite institutions, the correlation between listening to jazz and being hard{-}working is {-}0.09. If we look at students accepted to elite institutions, does it mean that listening to jazz affects being hard{-}working?} \\
\texttt{Start your answer with 'Yes' or 'No', followed by additional reasoning or evidence to support your explanation.} \\
            }
        }
    }
     \\ \\

    \theutterance \stepcounter{utterance}  
    & \multicolumn{4}{p{0.6\linewidth}}{
        \cellcolor[rgb]{0.561,0.737,0.561}{
            \makecell[{{p{\linewidth}}}]{
                \texttt{\tiny{[P1$\rangle$GM]}}
                \texttt{No.} \\
\\ 
\texttt{The correlation between listening to jazz and being hard{-}working among students accepted to elite institutions is {-}0.09, indicating a negative relationship. However, this does not imply causation from listening to jazz to being hard{-}working. } \\
\\ 
\texttt{In this hypothetical world, we have two direct effects on elite institution admission status: listening to jazz and effort. The correlation between listening to jazz and being hard{-}working among students accepted to elite institutions is likely due to the fact that both listening to jazz and being hard{-}working are correlated with a third variable, which is the effort that directly affects elite institution admission status.} \\
\\ 
\texttt{In other words, students who are hard{-}working are more likely to be accepted to elite institutions, and among these students, those who listen to jazz are likely to be a subset of those who are hard{-}working. This is an example of a confounding variable, where the relationship between listening to jazz and being hard{-}working is confounded by the underlying variable of effort.} \\
\\ 
\texttt{To establish causality from listening to jazz to being hard{-}working, we would need to control for the confounding variable of effort or show that the relationship between listening to jazz and being hard{-}working is not due to the effect of effort on both variables. However, based on the information provided, it is not possible to conclude that listening to jazz affects being hard{-}working.} \\
            }
        }
    }
    & & \\ \\

    \theutterance \stepcounter{utterance}  
    & & & \multicolumn{2}{p{0.3\linewidth}}{
        \cellcolor[rgb]{0.9,0.9,0.9}{
            \makecell[{{p{\linewidth}}}]{
                \texttt{\tiny{[GM$|$GM]}}
                \texttt{No} \\
            }
        }
    }
    & & \\ \\

    \theutterance \stepcounter{utterance}  
    & & & \multicolumn{2}{p{0.3\linewidth}}{
        \cellcolor[rgb]{0.9,0.9,0.9}{
            \makecell[{{p{\linewidth}}}]{
                \texttt{\tiny{[GM$|$GM]}}
                \texttt{no} \\
            }
        }
    }
    & & \\ \\

    \theutterance \stepcounter{utterance}  
    & & & \multicolumn{2}{p{0.3\linewidth}}{
        \cellcolor[rgb]{0.9,0.9,0.9}{
            \makecell[{{p{\linewidth}}}]{
                \texttt{\tiny{[GM$|$GM]}}
                \texttt{game\_result = WIN} \\
            }
        }
    }
    & & \\ \\

\end{supertabular}
}

\end{document}
